\section{对齐与治理：治理的对象是人，不是机器}

\subsection{对齐的悖论}

\begin{warningbox}{对齐到底对齐什么？}
机器一定要和人类对齐吗？人类自身存在贪婪、欺骗等缺陷——而这些恰恰是机器原本没有的。人类道德并非最高标准。那么``对齐''到底该对齐什么？这值得深入讨论。
\end{warningbox}

\subsection{治理研究者与使用者}

张钹院士明确指出：\textbf{最重要的治理对象不是机器，而是人类}——特别是研究者和使用者。AI时代的企业家必须承担更大的社会责任。


